\documentclass{../template/note}

\author{Oliver West}
\title{Imperial maths}
\date{}

\begin{document}
    
    \maketitle

    \tableofcontents

    \lecture{Lecture 1 -- 09/12/25}

    \section{Catenaries}

    Consider a chain of length $L$ and constant density per unit length $\rho_0$, hanging from fixtures $A$ and $B$ with $AB < L$. What will the shape of the chain be as it hangs?

    Let us make the centre of the chain, $P$, the origin. Then $B = (x_0, y_0)$.

    Pick an arbitrary point on the chain, $Q = (x, y)$, and consider the piece of chain $PQ$. Let us consider the forces acting:
    \begin{itemize}
        \item We have a tension force $\vec T(x)$, tangent to the chain at $Q$.
        \item We have a constant tension $T_0$ tangent to the chain at $P$.
        \item We have the weight due to gravity of $PQ$.
    \end{itemize}
    Horizontally, we have
    \begin{align}
        T_0 = T(x) \cos \theta
    \end{align}
    and vertically we have
    \begin{align}
        T(x) \sin \theta = PQ \rho_0 g
    \end{align}
    where $PQ$ is the length of the chain between $P$ and $Q$. Assuming the shape is given by $y(x)$, we can calculate
    \begin{align}
        PQ = \int_0^x \left(1 + y'^2\right)^{1/2} \; dx
    \end{align}
    Combining these equations, we have
    \begin{align}
        \rho_0 g \int_0^x \left(1 + y'^2\right)^{1/2} \; dt = T_0 \tan \theta = T_0 y'
    \end{align}
    since $\tan \theta = \frac{dy}{dx}$. To eliminate the integral, let us differentiate both sides to obtain
    \begin{align*}
        y'' = \frac{\rho_0g}{T_0}\left(1+y'^2\right)^{1/2}
    \end{align*}
    which is unfortunately non-linear. However, we might first wish to put $y' = p$, giving
    \begin{align*}
        \frac{dp}{dx} = \frac{\rho_0g}{T_0}\left(1+p^2\right)^{1/2}
    \end{align*}
    and this is now a separable equation! We have
    \begin{align*}
        \int \frac{1}{(1+p^2)^{1/2}} \; dp &= \int \frac{\rho_0g}{T_0} \; dx
    \end{align*}
    The right-hand side integral is simple, but for the left side, we need to use that
    \[
        \cosh^2 t - \sinh^2 t = 1 \iff \cosh^2 t = 1 + \sinh^2 t
    \]
    so we should substitute $p = \sinh t$, $dp = \cosh t \; dt$. We get
    \begin{align*}
        \int \frac{\cosh t}{(1+\sinh^2t)^{1/2}} \; dt &= \int 1 \; dt \\
        &= t = \arsinh p \\
        &= \arsinh\left(\frac{dy}{dx}\right)
    \end{align*}
    So we have
    \begin{align*}
        \arsinh\left(\frac{dy}{dx}\right) &= \frac{\rho_0g}{T_0}x + C
    \end{align*}
    for some constant $C$. In fact, since $y' = 0$ at $x=0$, $C = 0$. We have
    \begin{align*}
        \frac{dy}{dx} &= \sinh\left(\frac{\rho_0g}{T_0}x\right) \\
        \implies y &= \frac{T_0}{\rho_0g}\cosh\left(\frac{\rho_0g}{T_0}x\right) + D
    \end{align*}
    Since $x=0$, $y=0$, we get
    \[ 
        D = -\frac{T_0}{\rho_0g}
    \]
    and hence we conclude
    \[
        \boxed{y = \frac{T_0}{\rho_0g}\left(\cosh\left(\frac{\rho_0g}{T_0}x\right)-1\right)}
    \]
    In practice, we probably wish to do this in terms of the length. We have
    \begin{align*}
        \frac{L}{2} &= \int_0^{x_0} \left(1+y'^2\right)^{1/2} \; dx \\
        &= \frac{T_0}{\rho_0g}y' && \text{(by 4)}\\
        &= \frac{T_0}{\rho_0g}\sinh\left(\frac{\rho_0g}{T_0}x\right)
    \end{align*}
    so we can numerically solve for $T_0$ in terms of $L$, $\rho_0$ and $g$. 
    
    \section{Power series solutions of differential equations}

    Consider a simple differential equation:
    \[
        \frac{dy}{dx} = y
    \]
    with $y(0) = y_0$. We know the solution is given by $y = y_0e^x$. But let's try solving this another, more generalisable way. Suppose
    \begin{align*}
        y &= \sum_{n=0}^\infty a_n x^n \\
        \implies y' &= \sum_{n=1}^\infty na_n x^{n-1}
    \end{align*}
    Thus we have
    \begin{align*}
        \sum_{n=0}^\infty a_n x^n &= \sum_{n=1}^\infty na_n x^{n-1} \\
        &= \sum_{n=0}^\infty (n+1)a_{n+1}x^n
    \end{align*}
    Subtracting, we have
    \begin{align*}
        \implies \sum_{n=0}^\infty [a_n - (n+1)a_{n+1}]x^n &= 0 \\
        \implies a_n &= (n+1)a_{n+1}
    \end{align*}
    for all $n$. From this it becomes clear that
    \[
        a_n = \frac{a_0}{n!}
    \]
    so
    \begin{align*}
        y &= \sum_{n=0}^\infty \frac{a_0}{n!} x^n \\
        &= a_0 \sum_{n=0}^\infty \frac{1}{n!} x^n \\
        &= a_0 e^x
    \end{align*}
    and the initial condition gives $a_0 = y_0$ as required.
    
    
    
    

\end{document}